\chapter*{\textarabic{ملخص}}
\addcontentsline{toc}{chapter}{\textarabic{ملخص}}

\begin{otherlanguage}{arabic}
تُستعمل قوائم مكونات البرمجيات، أو \textenglish{SBOM}، بشكل متزايد لدعم
شفافية سلسلة توريد البرمجيات وإدارة الثغرات. غير أن وجود \textenglish{SBOM}
صحيحة من الناحية التركيبية لا يعني بالضرورة أنها موثوقة، خاصة عندما يعاد
بناؤها من تطبيق أندرويد مترجم دون الوصول إلى الشفرة المصدرية أو ملفات بيان
الاعتماديات. فقد تؤدي الترجمة، والتحسين، والتعمية، وأسماء عائلات القطع
البرمجية، ونقص قواعد المراجع، والتشابه بين الإصدارات إلى إضعاف الأدلة التي
تدعم ادعاء وجود مكون معين أو إلى مناقضتها.

يدرس هذا العمل استرجاع \textenglish{SBOM} الموثوقة بوصفه مسألة إسناد
مكونات برمجية قائمة على الأدلة. وينظم الأبحاث السابقة حول كشف المكتبات
الخارجية في أندرويد واسترجاع المكونات الثنائية ضمن تصنيفات للتقنيات، ومصادر
الأدلة، ومخاطر التحويل، وممارسات التقييم. وبناء على الفجوات الناتجة، يقترح
إطارا مفهوميا يسجل لكل ادعاء مسترجع هوية المكون، والإصدار أو مجموعة
الإصدارات، ومصدر الأدلة، وجودتها، ومخاطر التحويل، ومستوى الثقة، والغموض،
وفئة التحقق، والتفسير الأمني المسموح به لاحقا.

يُستخدم نظام مشروع نهاية الدراسة لتحليل \textenglish{APK} وإنتاج
\textenglish{SBOM} كدراسة حالة محدودة. وتبين أدلته المضبوطة، والنتائج
النظيفة والجزئية المحققة بالمصدر، والحدود العليا المعروفة بالمصدر، والنتائج
التجريبية أو الإرشادية أو الحدية، ضرورة الفصل بين هوية المكتبة، والإصدار
الدقيق، وتغطية المراجع، واسترجاع المطابق، وسلطة ادعاءات الثغرات. ويؤكد تقييم
دراسة الحالة أن الإطار يحافظ على مسارات الأدلة هذه ويمنع رفع سلطة الادعاء
عندما تفشل التغييرات التي تبدو مفيدة في اختبارات الغموض أو الضبط أو الترقية.
لذلك يفضل
الإطار المقترح معرفة جزئية قابلة للتدقيق وامتناعا صريحا عن الحكم بدلا من
يقين غير مدعوم، كما يحدد برنامجا بحثيا لدمج الأدلة بصورة معايرة، ولبناء
\textenglish{SBOM} واعية بالمصدر، ولاختبارات واقعية، ولاسترجاع المكونات
الأصلية أو المحملة ديناميكيا.
\end{otherlanguage}

\vspace{0.5cm}
\noindent\textbf{\textarabic{الكلمات المفتاحية} ---}
\begin{otherlanguage}{arabic}
قائمة مكونات البرمجيات، سلسلة توريد البرمجيات، أندرويد، إسناد المكونات،
مصدر الأدلة، دمج الأدلة، الثقة، عدم اليقين، إدارة الثغرات.
\end{otherlanguage}
